\documentclass[12pt]{article}

% Separate working manuscript. Canonical main.* files are unchanged.
% !TEX TS-program = xelatex
% File: preamble.tex
% Purpose: Brett Reynolds house style LaTeX preamble
% Version: 2.0.0 (December 2025 typography redesign)
%
% Usage: % !TEX TS-program = xelatex
% File: preamble.tex
% Purpose: Brett Reynolds house style LaTeX preamble
% Version: 2.0.0 (December 2025 typography redesign)
%
% Usage: % !TEX TS-program = xelatex
% File: preamble.tex
% Purpose: Brett Reynolds house style LaTeX preamble
% Version: 2.0.0 (December 2025 typography redesign)
%
% Usage: \input{.house-style/preamble.tex} in main document
%
% Compiler: XeLaTeX (default) or LuaLaTeX
%           NOT pdfLaTeX - requires fontspec for fonts

% ===========================
% FONTS
% ===========================
\usepackage{fontspec}

\setmainfont{EBGaramond}[
  Extension=.otf,
  UprightFont=*-Regular,
  ItalicFont=*-Italic,
  BoldFont=*-Bold,
  BoldItalicFont=*-BoldItalic,
  Numbers=OldStyle,
  Ligatures=TeX,
]

% IPA fallback font
\newfontfamily\ipafont{CharisSIL}[
  Extension=.ttf,
  UprightFont=*-Regular,
  ItalicFont=*-Italic,
  BoldFont=*-Bold,
  BoldItalicFont=*-BoldItalic,
]
\newcommand{\ipa}[1]{{\ipafont #1}}

% Lining figures when needed (tables, years in isolation)
\providecommand{\liningnums}[1]{{\addfontfeatures{Numbers=Lining}#1}}

% Monospace
\setmonofont{Inconsolatazi4}[
  Extension=.otf,
  UprightFont=*-Regular,
  BoldFont=*-Bold,
  Scale=MatchLowercase,
]

% ===========================
% PAGE LAYOUT
% ===========================
\usepackage[
  letterpaper,
  inner=1.25in,
  outer=1in,
  top=1in,
  bottom=1.25in,
  marginparwidth=0.6in,
]{geometry}

\usepackage[british]{babel}
\usepackage[final]{microtype}

% ===========================
% HEADINGS
% ===========================
\usepackage{titlesec}

% Section: small caps, number in margin
\titleformat{\section}
  {\normalfont\scshape}
  {\llap{\thesection\quad}}
  {0pt}
  {}

% Subsection: small caps sentence case
\titleformat{\subsection}
  {\normalfont\scshape}
  {\thesubsection\quad}
  {0pt}
  {}

% Subsubsection: upright
\titleformat{\subsubsection}
  {\normalfont}
  {\thesubsubsection\quad}
  {0pt}
  {}

% Spacing
\titlespacing*{\section}{0pt}{2ex plus 1ex minus .2ex}{1ex plus .2ex}
\titlespacing*{\subsection}{0pt}{1.5ex plus 1ex minus .2ex}{0.5ex plus .2ex}
\titlespacing*{\subsubsection}{0pt}{1ex plus 0.5ex minus .1ex}{0.3ex plus .1ex}

% ===========================
% RUNNING HEADS
% ===========================
\usepackage{fancyhdr}
\pagestyle{fancy}
\fancyhf{}
\setlength{\headheight}{13.6pt}
\renewcommand{\headrulewidth}{0pt}

% For oneside documents (most papers):
\fancyhead[L]{\small\scshape\leftmark}
\fancyhead[R]{\small\thepage}
\fancyfoot{}

% ===========================
% COLORS & HYPERLINKS
% ===========================
\usepackage{xcolor}
\definecolor{linkmaroon}{RGB}{128, 0, 32}

\usepackage{hyperref}
\hypersetup{
  colorlinks=true,
  linkcolor=linkmaroon,
  citecolor=linkmaroon,
  urlcolor=linkmaroon,
  pdfauthor={Brett Reynolds},
}

\usepackage{orcidlink}

% ===========================
% QUOTATIONS
% ===========================
\usepackage[style=american]{csquotes} % Double quotes outer (Oxford spelling kept via babel)
\MakeOuterQuote{"}

% ===========================
% SEMANTIC MACROS
% ===========================
% Terms (linguistic concepts being introduced/defined)
\newcommand{\term}[1]{\textsc{#1}}

% Mentions (words/expressions under discussion)
\newcommand{\mention}[1]{\textit{#1}}

% Mentions in headings (angle brackets instead of italics)
% Robust: headings are moving arguments (written to .toc/.aux)
\DeclareRobustCommand{\mentionhead}[1]{⟨\textup{#1}⟩}

% Object language (cited forms, foreign words)
\newcommand{\olang}[1]{\textit{#1}}

% Small-caps abbreviations for glosses
\newcommand{\abbr}[1]{\textsc{#1}}

% Cross-linguistic subscript marker
\usepackage{marvosym}
\newcommand{\crossmark}{\textsubscript{\Cross}}

% ===========================
% MATHS AND SYMBOLS
% ===========================
\usepackage{amsmath,amssymb}

% ===========================
% LINGUISTIC EXAMPLES
% ===========================
\usepackage{langsci-gb4e}
\makeatletter
\@ifundefined{noautomath}{}{\noautomath}
\makeatother

% Judgement markers
\newcommand{\ungram}[1]{*\!#1}
\newcommand{\marg}[1]{?\!#1}
\newcommand{\odd}[1]{\#\!#1}

% ===========================
% LISTS
% ===========================
\usepackage{enumitem}
\setlist{nosep, leftmargin=*}
\setlist[enumerate]{label=\arabic*.}
\setlist[itemize]{label=--}

% ===========================
% TABLES & FIGURES
% ===========================
\usepackage{booktabs}
\usepackage{array}
\usepackage{graphicx}
\graphicspath{{figures/}}

% ===========================
% BIBLIOGRAPHY
% ===========================
\usepackage[backend=biber,style=apa,natbib=true,doi=true,isbn=false,url=true]{biblatex}
\addbibresource{references.bib}
\IfFileExists{references-local.bib}{\addbibresource{references-local.bib}}{}

% Possessive citation: Author's (Year)
\newcommand{\posscite}[1]{\citeauthor{#1}'s (\citeyear{#1})}

% ===========================
% UTILITIES
% ===========================
\usepackage{xspace}
\newcommand{\eg}{e.g.,\xspace}
\newcommand{\ie}{i.e.,\xspace}
\newcommand{\etc}{etc.\xspace}
 in main document
%
% Compiler: XeLaTeX (default) or LuaLaTeX
%           NOT pdfLaTeX - requires fontspec for fonts

% ===========================
% FONTS
% ===========================
\usepackage{fontspec}

\setmainfont{EBGaramond}[
  Extension=.otf,
  UprightFont=*-Regular,
  ItalicFont=*-Italic,
  BoldFont=*-Bold,
  BoldItalicFont=*-BoldItalic,
  Numbers=OldStyle,
  Ligatures=TeX,
]

% IPA fallback font
\newfontfamily\ipafont{CharisSIL}[
  Extension=.ttf,
  UprightFont=*-Regular,
  ItalicFont=*-Italic,
  BoldFont=*-Bold,
  BoldItalicFont=*-BoldItalic,
]
\newcommand{\ipa}[1]{{\ipafont #1}}

% Lining figures when needed (tables, years in isolation)
\providecommand{\liningnums}[1]{{\addfontfeatures{Numbers=Lining}#1}}

% Monospace
\setmonofont{Inconsolatazi4}[
  Extension=.otf,
  UprightFont=*-Regular,
  BoldFont=*-Bold,
  Scale=MatchLowercase,
]

% ===========================
% PAGE LAYOUT
% ===========================
\usepackage[
  letterpaper,
  inner=1.25in,
  outer=1in,
  top=1in,
  bottom=1.25in,
  marginparwidth=0.6in,
]{geometry}

\usepackage[british]{babel}
\usepackage[final]{microtype}

% ===========================
% HEADINGS
% ===========================
\usepackage{titlesec}

% Section: small caps, number in margin
\titleformat{\section}
  {\normalfont\scshape}
  {\llap{\thesection\quad}}
  {0pt}
  {}

% Subsection: small caps sentence case
\titleformat{\subsection}
  {\normalfont\scshape}
  {\thesubsection\quad}
  {0pt}
  {}

% Subsubsection: upright
\titleformat{\subsubsection}
  {\normalfont}
  {\thesubsubsection\quad}
  {0pt}
  {}

% Spacing
\titlespacing*{\section}{0pt}{2ex plus 1ex minus .2ex}{1ex plus .2ex}
\titlespacing*{\subsection}{0pt}{1.5ex plus 1ex minus .2ex}{0.5ex plus .2ex}
\titlespacing*{\subsubsection}{0pt}{1ex plus 0.5ex minus .1ex}{0.3ex plus .1ex}

% ===========================
% RUNNING HEADS
% ===========================
\usepackage{fancyhdr}
\pagestyle{fancy}
\fancyhf{}
\setlength{\headheight}{13.6pt}
\renewcommand{\headrulewidth}{0pt}

% For oneside documents (most papers):
\fancyhead[L]{\small\scshape\leftmark}
\fancyhead[R]{\small\thepage}
\fancyfoot{}

% ===========================
% COLORS & HYPERLINKS
% ===========================
\usepackage{xcolor}
\definecolor{linkmaroon}{RGB}{128, 0, 32}

\usepackage{hyperref}
\hypersetup{
  colorlinks=true,
  linkcolor=linkmaroon,
  citecolor=linkmaroon,
  urlcolor=linkmaroon,
  pdfauthor={Brett Reynolds},
}

\usepackage{orcidlink}

% ===========================
% QUOTATIONS
% ===========================
\usepackage[style=american]{csquotes} % Double quotes outer (Oxford spelling kept via babel)
\MakeOuterQuote{"}

% ===========================
% SEMANTIC MACROS
% ===========================
% Terms (linguistic concepts being introduced/defined)
\newcommand{\term}[1]{\textsc{#1}}

% Mentions (words/expressions under discussion)
\newcommand{\mention}[1]{\textit{#1}}

% Mentions in headings (angle brackets instead of italics)
% Robust: headings are moving arguments (written to .toc/.aux)
\DeclareRobustCommand{\mentionhead}[1]{⟨\textup{#1}⟩}

% Object language (cited forms, foreign words)
\newcommand{\olang}[1]{\textit{#1}}

% Small-caps abbreviations for glosses
\newcommand{\abbr}[1]{\textsc{#1}}

% Cross-linguistic subscript marker
\usepackage{marvosym}
\newcommand{\crossmark}{\textsubscript{\Cross}}

% ===========================
% MATHS AND SYMBOLS
% ===========================
\usepackage{amsmath,amssymb}

% ===========================
% LINGUISTIC EXAMPLES
% ===========================
\usepackage{langsci-gb4e}
\makeatletter
\@ifundefined{noautomath}{}{\noautomath}
\makeatother

% Judgement markers
\newcommand{\ungram}[1]{*\!#1}
\newcommand{\marg}[1]{?\!#1}
\newcommand{\odd}[1]{\#\!#1}

% ===========================
% LISTS
% ===========================
\usepackage{enumitem}
\setlist{nosep, leftmargin=*}
\setlist[enumerate]{label=\arabic*.}
\setlist[itemize]{label=--}

% ===========================
% TABLES & FIGURES
% ===========================
\usepackage{booktabs}
\usepackage{array}
\usepackage{graphicx}
\graphicspath{{figures/}}

% ===========================
% BIBLIOGRAPHY
% ===========================
\usepackage[backend=biber,style=apa,natbib=true,doi=true,isbn=false,url=true]{biblatex}
\addbibresource{references.bib}
\IfFileExists{references-local.bib}{\addbibresource{references-local.bib}}{}

% Possessive citation: Author's (Year)
\newcommand{\posscite}[1]{\citeauthor{#1}'s (\citeyear{#1})}

% ===========================
% UTILITIES
% ===========================
\usepackage{xspace}
\newcommand{\eg}{e.g.,\xspace}
\newcommand{\ie}{i.e.,\xspace}
\newcommand{\etc}{etc.\xspace}
 in main document
%
% Compiler: XeLaTeX (default) or LuaLaTeX
%           NOT pdfLaTeX - requires fontspec for fonts

% ===========================
% FONTS
% ===========================
\usepackage{fontspec}

\setmainfont{EBGaramond}[
  Extension=.otf,
  UprightFont=*-Regular,
  ItalicFont=*-Italic,
  BoldFont=*-Bold,
  BoldItalicFont=*-BoldItalic,
  Numbers=OldStyle,
  Ligatures=TeX,
]

% IPA fallback font
\newfontfamily\ipafont{CharisSIL}[
  Extension=.ttf,
  UprightFont=*-Regular,
  ItalicFont=*-Italic,
  BoldFont=*-Bold,
  BoldItalicFont=*-BoldItalic,
]
\newcommand{\ipa}[1]{{\ipafont #1}}

% Lining figures when needed (tables, years in isolation)
\providecommand{\liningnums}[1]{{\addfontfeatures{Numbers=Lining}#1}}

% Monospace
\setmonofont{Inconsolatazi4}[
  Extension=.otf,
  UprightFont=*-Regular,
  BoldFont=*-Bold,
  Scale=MatchLowercase,
]

% ===========================
% PAGE LAYOUT
% ===========================
\usepackage[
  letterpaper,
  inner=1.25in,
  outer=1in,
  top=1in,
  bottom=1.25in,
  marginparwidth=0.6in,
]{geometry}

\usepackage[british]{babel}
\usepackage[final]{microtype}

% ===========================
% HEADINGS
% ===========================
\usepackage{titlesec}

% Section: small caps, number in margin
\titleformat{\section}
  {\normalfont\scshape}
  {\llap{\thesection\quad}}
  {0pt}
  {}

% Subsection: small caps sentence case
\titleformat{\subsection}
  {\normalfont\scshape}
  {\thesubsection\quad}
  {0pt}
  {}

% Subsubsection: upright
\titleformat{\subsubsection}
  {\normalfont}
  {\thesubsubsection\quad}
  {0pt}
  {}

% Spacing
\titlespacing*{\section}{0pt}{2ex plus 1ex minus .2ex}{1ex plus .2ex}
\titlespacing*{\subsection}{0pt}{1.5ex plus 1ex minus .2ex}{0.5ex plus .2ex}
\titlespacing*{\subsubsection}{0pt}{1ex plus 0.5ex minus .1ex}{0.3ex plus .1ex}

% ===========================
% RUNNING HEADS
% ===========================
\usepackage{fancyhdr}
\pagestyle{fancy}
\fancyhf{}
\setlength{\headheight}{13.6pt}
\renewcommand{\headrulewidth}{0pt}

% For oneside documents (most papers):
\fancyhead[L]{\small\scshape\leftmark}
\fancyhead[R]{\small\thepage}
\fancyfoot{}

% ===========================
% COLORS & HYPERLINKS
% ===========================
\usepackage{xcolor}
\definecolor{linkmaroon}{RGB}{128, 0, 32}

\usepackage{hyperref}
\hypersetup{
  colorlinks=true,
  linkcolor=linkmaroon,
  citecolor=linkmaroon,
  urlcolor=linkmaroon,
  pdfauthor={Brett Reynolds},
}

\usepackage{orcidlink}

% ===========================
% QUOTATIONS
% ===========================
\usepackage[style=american]{csquotes} % Double quotes outer (Oxford spelling kept via babel)
\MakeOuterQuote{"}

% ===========================
% SEMANTIC MACROS
% ===========================
% Terms (linguistic concepts being introduced/defined)
\newcommand{\term}[1]{\textsc{#1}}

% Mentions (words/expressions under discussion)
\newcommand{\mention}[1]{\textit{#1}}

% Mentions in headings (angle brackets instead of italics)
% Robust: headings are moving arguments (written to .toc/.aux)
\DeclareRobustCommand{\mentionhead}[1]{⟨\textup{#1}⟩}

% Object language (cited forms, foreign words)
\newcommand{\olang}[1]{\textit{#1}}

% Small-caps abbreviations for glosses
\newcommand{\abbr}[1]{\textsc{#1}}

% Cross-linguistic subscript marker
\usepackage{marvosym}
\newcommand{\crossmark}{\textsubscript{\Cross}}

% ===========================
% MATHS AND SYMBOLS
% ===========================
\usepackage{amsmath,amssymb}

% ===========================
% LINGUISTIC EXAMPLES
% ===========================
\usepackage{langsci-gb4e}
\makeatletter
\@ifundefined{noautomath}{}{\noautomath}
\makeatother

% Judgement markers
\newcommand{\ungram}[1]{*\!#1}
\newcommand{\marg}[1]{?\!#1}
\newcommand{\odd}[1]{\#\!#1}

% ===========================
% LISTS
% ===========================
\usepackage{enumitem}
\setlist{nosep, leftmargin=*}
\setlist[enumerate]{label=\arabic*.}
\setlist[itemize]{label=--}

% ===========================
% TABLES & FIGURES
% ===========================
\usepackage{booktabs}
\usepackage{array}
\usepackage{graphicx}
\graphicspath{{figures/}}

% ===========================
% BIBLIOGRAPHY
% ===========================
\usepackage[backend=biber,style=apa,natbib=true,doi=true,isbn=false,url=true]{biblatex}
\addbibresource{references.bib}
\IfFileExists{references-local.bib}{\addbibresource{references-local.bib}}{}

% Possessive citation: Author's (Year)
\newcommand{\posscite}[1]{\citeauthor{#1}'s (\citeyear{#1})}

% ===========================
% UTILITIES
% ===========================
\usepackage{xspace}
\newcommand{\eg}{e.g.,\xspace}
\newcommand{\ie}{i.e.,\xspace}
\newcommand{\etc}{etc.\xspace}

\usepackage{xurl}
\usepackage[normalem]{ulem}
\setlength{\ULdepth}{1.8pt}
\setlength{\emergencystretch}{3em}
\setlength{\headheight}{14pt}
\newcommand{\exmention}[1]{\mention{#1}}
\newcommand{\exfocus}[1]{\uline{#1}}
\addbibresource{inquiry-relative-categories.bib}

% Publisher-verified correction without changing the shared bibliography.
\DeclareSourcemap{
  \maps[datatype=bibtex]{
    \map[overwrite=true]{
      \step[fieldsource=entrykey, match=\regexp{\Aunuabonahdaniel2020\Z}, final]
      \step[fieldset=doi, fieldvalue={10.1016/j.pragma.2020.04.004}]
    }
  }
}

\newif\ifblind
\blindfalse
\ifblind
\DeclareSourcemap{
  \maps[datatype=bibtex]{
    \map[overwrite=true]{
      \step[fieldsource=entrykey, match=\regexp{\Areynolds}, final]
      \step[fieldset=author, fieldvalue={Author}]
      \step[fieldset=title, fieldvalue={[Title removed for anonymous review]}]
      \step[fieldset=note, null]
      \step[fieldset=url, null]
      \step[fieldset=doi, null]
    }
  }
}
\fi
\hypersetup{
  pdftitle={Inquiry-Relative Categories: English Interjections as a Worked Case},
  pdfkeywords={inquiry-relative categories, category pluralism, projectibility, interjections, sign-types, path dependence},
}
\ifblind
\hypersetup{pdfauthor={}}
\else
\hypersetup{pdfauthor={Brett Reynolds}}
\fi
\title{Inquiry-Relative Categories:\\English Interjections as a Worked Case}
\ifblind
\author{[Author removed for blind review]}
\else
\author{Brett Reynolds \orcidlink{0000-0003-0073-7195}%
\thanks{Contact: \href{mailto:brett.reynolds@humber.ca}{brett.reynolds@humber.ca}}\\
Humber Polytechnic \& University of Toronto}
\fi
\date{\today}

\begin{document}
\maketitle

\begin{abstract}
Linguistic category disputes often seek one classification for expressions
that support several kinds of inference. I develop an account of
inquiry-relative categories through English interjections. Established
grammatical descriptions support a core including \mention{ouch},
\mention{oops}, \mention{hey}, and \mention{sh}: their conventional uses
connect independent utterance with clausal supplementation and restricted
integration. This profile licenses grammatical expectations across
environments while grouping expressions with different meanings and
conversational jobs. Expressive contribution, in turn, connects standalone
\mention{Damn!} with prenominal \mention{damn} across a grammatical
difference. Comparing stance with sound depiction in \mention{Bang!}
shows why even one semantic inquiry can require several distinctions.
These cases support a division of analytical labour: grammatical
membership, expressive commonality, and depictive contribution answer
different questions about overlapping material. The account distinguishes
descriptive success from further projective and causal claims, whose
evidence requires separate assessment. Several classifications can
be correct without one superseding the others, while incompatible claims
about the same use remain open to empirical adjudication.
\end{abstract}

\noindent\textbf{Keywords:} inquiry-relative categories; category pluralism;
projectibility; interjections; sign-types; path dependence

\aidisclosure{Claude Opus 4.8, Claude Opus 5, Claude Opus 5.5, GPT-5.6 Sol,
GPT-6 Astra Pro, and GPT-6 (Codex)}

\section{Introduction}\label{sec:intro}

Standalone \mention{Damn!}, complement-taking \mention{Damn these mosquitoes!},
and prenominal \mention{damn} in \mention{the damn car} present different
combinations of grammatical and semantic properties. The modifier participates
in a noun phrase while retaining an expressive contribution. The standalone
form can constitute a stance display. The complement-taking expression has
internal structure even when the whole expression stands apart from a
surrounding clause. A single judgment about whether \mention{damn} is an
interjection leaves these differences unresolved.

A grammatical answer begins with a category whose members support
expectations about their use in constructions. \mention{Ouch!} can
constitute an utterance and supplement a clause in \mention{Ouch, that
hurt}. Attention-getting \mention{hey} and silencing \mention{sh} share
the broad grammatical profile despite their different conversational
jobs. Establishing that profile gives us a basis for assessing the
\mention{damn} uses without requiring all expressive signs to belong
to the same lexical category.

A comparison with depictive \mention{Bang!} exposes a further distinction.
The selected \mention{damn} use conveys an evaluative stance; the selected
\mention{Bang!} depicts a sound. Both concern meaning, but they motivate
different groupings and different questions about interpretation. Grouping
one under semantics and the other under semiotics would do nothing to
resolve those questions. Nor would calling both interjections establish
that they share the relevant semantic properties.

This paper asks what warrants recognizing a category for an inquiry.
An inquiry identifies a question and a standard of success; evidence
constrains which grouping answers it. A field can pursue several inquiries,
an inquiry can employ several categories, and a category can serve questions
pursued across fields. English interjections make this many-to-many relation
visible because grammatical, semantic, interactional, phonological, and
historical descriptions repeatedly organize overlapping material.

A grouping can make measurement comparable or description informative.
The distinctions among \mention{damn} uses have descriptive value even
before a new prediction is tested. \term{Projectibility} concerns the
warrant for expectations beyond the evidence used to identify a case
\citep{Goodman1955}. Neither descriptive usefulness nor a successful
projection identifies its own causal explanation. These achievements
need separate assessment.

Pluralism about useful and objectively grounded groupings has substantial
precedents. \textcite[5--7, 17--18]{dupre1993disorder} relates the choice of
classification to the goals of inquiry while rejecting one uniquely
privileged division of things. \textcite[63--64, 69--72]{khalidi2013}
defends crosscutting kinds answering different epistemic interests.
Their wider metaphysical commitments needn't all be adopted to recognize
that distinct questions can require overlapping groupings.

Within linguistics, \textcite[663, 666, 680]{haspelmath2010} distinguishes
language-particular descriptive categories from concepts designed for
cross-language comparison. The division of labour here is different:
several inquiries within English can organize the same expressions
differently. The assessment developed here starts with a positive grammatical
case and follows its relation to other groupings. A grammatical category
organizes expectations about constructions; an expressive grouping preserves
a relation that crosses its boundary. Each succeeds for a different
question. Specifying the bearer, grouping, and intended inference also
shows when two analyses compete and what evidence could decide between them.

The worked cases are named for the relations they investigate: grammatical
interjections, expressive signs, depictive signs, continuers, repair
initiators, and closing formats. They needn't all operate at the same level
of generality or retain an interjection label. In these comparisons, the
common bearer is a conventionalized sign-type, an identifiable expression
in a conventional use. Classifications of realizations or historical
changes require explicit relations to these sign-types.

The historical case concerns \term{path dependence}: development can leave
source-specific traces after uses converge. Gehweiler's history of
\mention{gee} supplies a documented explanation of one outcome; a verbal
route to complement-taking \mention{damn} remains partly reconstructive.
The corpus material has a narrower role, showing how tags and strings
can fail to identify the intended bearers. A linked supplement preserves
the detailed measurement checks and two withdrawn confirmatory analyses.

Section~\ref{sec:worked} establishes a grammatical core, traces expressive
commonality across grammatical differences, and compares expressive and depictive
contributions. Section~\ref{sec:assessment} distinguishes the achievements
of these groupings. Sections~\ref{sec:phonology} and~\ref{sec:history}
apply the assessment to phonological and historical claims;
Section~\ref{sec:extensions} considers broader social and processing
classes. The measurement case and discussion then draw out what changes
when the object, grouping, and claim are kept explicit.

\section{Several groupings of conventional signs}\label{sec:worked}

Grammars and studies of meaning and interaction organize the traditional
interjection domain differently. \textcite[90]{jespersen1924} treated
interjection as a manner in which words of other categories are used.
\textit{The Cambridge Grammar of the English Language} (\textit{CGEL})
recognizes a lexical category \citep[22]{HuddlestonPullum2002}. The
\textit{Longman Grammar of Spoken and Written English} places interjections
within the broader class of inserts, while assigning \mention{damn},
\mention{gosh}, and \mention{gee} to expletives
\citep[1082--1083, 1094--1095]{Biber1999}.

The disputes concern a well-studied domain. \textcite{libert2019} questions
the description of interjections as neglected; \textcite[63]{cram2008}
identifies their recurring role in controversies over the syntax--pragmatics
and semantics--pragmatics boundaries. The following comparisons preserve
what those traditions distinguish without treating each tradition as the
source of exactly one category.

\subsection{A grammatical core and what it projects}\label{sec:grammatical-core}

Someone drops a cup and says \mention{Oops!}; someone catches a finger
in a drawer and says \mention{Ouch!}. Either reaction can also accompany
a clause. The following constructed pairs illustrate the two environments.

\ea\label{ex:core-distribution}
  \ea \exmention{Oops!}\qquad\exmention{Oops, I've dropped it!}
  \ex \exmention{Ouch!}\qquad\exmention{Ouch, that hurt!}
  \z
\z

In the independent utterances, the lexical item supplies the expression
without an unspoken clause. In the longer examples, it functions as
a supplement to the clause. \mention{I've dropped it} and \mention{that
hurt} are already complete: the reaction fills no argument or integrated
modifier position within them. These are two characteristic environments
for the grammatical category \term{interjection}.

I take \mention{ah}, \mention{hey}, \mention{oh}, \mention{oops},
\mention{ouch}, \mention{sh}, \mention{ugh}, and \mention{wow}, in
their conventional interjectional uses, as a working core. \textit{CGEL}
identifies these as central members and attributes both independent
and supplementary use to them \citep[1361]{HuddlestonPullum2002}.
Its additional example \mention{alas} illustrates how a dated form
can have a secure grammatical assignment. Here, \term{core} concerns
the support for those assignments, rather than frequency or typicality.

The shared profile combines these positive possibilities with restricted
integration. The selected uses are characteristically invariant and don't
participate in the ordinary argument, predicative, and modifier patterns
of nouns, verbs, and adjectives. \textcite[105--106]{ameka1992} describes
primary interjections through their independent utterance capacity, limited
combination, and lack of productive inflection. \textcite[124--125]{wilkins1992}
likewise identifies a formal profile of conventional lexical utterances,
keeping their semantic and pragmatic properties available for separate
investigation.

This category licenses an inference between environments. An analyst who
has established the grammatical assignment of reaction \mention{ouch}
from independent uses can expect it also to supplement an intact
clause. From an identified supplementary \mention{oops}, the analyst can
expect an independent utterance without recovering a missing predicate.
The grammatical descriptions warrant these expectations about the core;
the pairs in (\ref{ex:core-distribution}) illustrate their application.
Neither the examples nor the category claim imply free combination
with every clause or a measured rate of predictive success.

An isolated word can have a different analysis. Ordinary \mention{Stop!}
is an imperative clause: the verb retains the possibilities illustrated
by \mention{Stop the car!} and \mention{Don't stop!}. A reply
\mention{Coffee} to \mention{What would you like?} retains nominal
organization, also visible in \mention{Some coffee}. These constructed
comparisons extend the contrast drawn by \textcite[124--125]{wilkins1992}
with adjectival \mention{Black} as an answer to a question about colour.
Utterance status alone doesn't reclassify the word.

Supplement function is shared by several categories. Compare the interjection
with the relative clause and noun phrase in these examples, adapted from
\textit{CGEL} \citep[1356--1357, 1360]{HuddlestonPullum2002}.

\ea\label{ex:supplements}
  \ea \exmention{\exfocus{Ah}, so you were there after all!}
  \ex \exmention{We called in to see Sue's parents, \exfocus{which made us rather late}.}
  \ex \exmention{A university lecturer, \exfocus{Dr Brown}, was arrested for the crime.}
  \z
\z

The relative clause and noun phrase retain their internal structure
and other distributional possibilities. The interjection's assignment rests
on the broader lexical profile. Treating all three as supplements
answers a question about their relation to an anchor; distinguishing
their categories answers questions about their internal organization and
distribution elsewhere. The two classifications can both be correct.

That grammatical core accommodates different contributions. A speaker can
use \mention{Hey!} to attract attention and \mention{Sh!} to request
silence, alongside \mention{Ouch!} as a reaction to pain.
\textcite[113, 115]{ameka1992} groups the first two as conative,
directed toward a listener's attention or action. Their shared grammar
doesn't identify a single action or require the evaluative or affective
stance characteristic of the expressive grouping developed below.

\textit{CGEL}'s broad characterization of interjection meaning as expressive
\citep[1361]{HuddlestonPullum2002} is compatible with recognizing these
differences. It supplies no reason to identify all its grammatical
members with a narrower semantic grouping. \textcite[114]{ameka1992}
also allows functional overlap: an attention-getter can convey affect.
The grammatical category is useful because of its distributional
expectations, whether or not a particular use also belongs to an
expressive or interactional grouping.\footnote{The grammatical analysis can
itself be disputed. \textcite[126--127 n.~5]{wilkins1992} treats \mention{oh} in \mention{Oh, don't go
on like an old fool} as a particle use, while reporting Ameka's
co-utterance analysis. The core claim
concerns identified conventional uses, not every occurrence of a spelling.}

\subsection{Expressive commonality across a grammatical boundary}\label{sec:damn-profiles}

Standalone \mention{Damn!} and supplementary \mention{damn} in
\mention{Damn, we're going to be late!} extend the grammatical comparison
to a form with an independently available verbal use. \textit{CGEL}
analyzes the exclamatory uses as reanalysed interjections
\citep[1361]{HuddlestonPullum2002}. Their distribution supplies the link
to the core; their conveyed stance supplies a different link to
expressive uses elsewhere in the grammar.

Prenominal \mention{damn}, in \mention{the damn car}, occurs inside the
noun phrase while contributing an attitude toward its referent.
\textcite[171, 182]{potts2007} distinguishes that expressive contribution
from descriptive content. Grammatical integration leaves a
semantic relation to standalone \mention{Damn!} available for analysis.
The grammatical analysis separates these uses; the expressive analysis
groups them together. Each preserves a relation needed for its own
question.

Complement-taking \mention{Damn these mosquitoes!} presents a further
grammatical question. The phrase's external detachment is compatible with
complementation inside it. \textit{CGEL} tentatively treats the head as
an exceptional interjection taking a noun-phrase complement
\citep[1361 n.~70]{HuddlestonPullum2002}. A verbal analysis remains a
substantive alternative. This disagreement concerns the head's grammar;
recognizing expressive commonality with the other \mention{damn} uses
leaves it intact.

The semantic grouping here is \term{expressive signs}: conventional uses
whose deployment conveys an evaluative or affective stance, rather than
merely describing that stance. This intended contribution is narrower than
all non-descriptive or use-conditional meaning. A conventional instruction
for managing interpretation or interaction needn't express such a stance.
Potts's diagnostics include perspective dependence, immediacy, and
intensifying repetition \citep[166--167]{potts2007}. For \mention{damn},
these ask whose attitude is conveyed, whether the utterance displays
that attitude on the occasion of speaking, and whether repeating the
form strengthens it. Whether those properties travel together across
the proposed extension is an empirical question.

Expressive commonality extends across the grammatical boundary. Prenominal
\mention{damn} qualifies because of its conveyed stance, while the
grammatical core includes conative uses for which that stance isn't
required. A fuller account of expressives can include integrated signs
with no traditional interjection label. Selecting examples from interjection
disputes lets us compare the two groupings without adding a semantic
membership condition based on that traditional label.

\subsection{Expressive and depictive signs}\label{sec:expressive-depictive}

Compare the selected expressive use of \mention{Damn!} with a sound depiction
\mention{Bang!} in a narrative. The bearer in each case is the conventional
use, not every occurrence of the spelling. The question is what contribution
the use makes to interpretation. The strongest simple alternative groups
both as expressive interjections on the basis of their vividness and
standalone packaging. That alternative still has to explain the difference
between conveying an attitude and presenting a sound.

A \term{depictive sign} conventionally presents a sound, event, or sensory
scene through a resemblance between aspects of its performance and what
it depicts. Dingemanse's account of ideophones supplies the wider comparison:
conventional lexical items can depict sensory imagery, and their grammatical
organization varies across languages \citep[2--3]{dingemanse2018}.
The term \term{expressive} is itself used for ideophones in some traditions;
that usage doesn't identify them with Potts's expressives.

For the selected \mention{Bang!}, the relevant contribution is a portrayal
of an impact sound. For \mention{Damn!}, it's the conveyed stance toward
a situation. The distinction changes the interpretation of a possible
modification. In a constructed narrative, \mention{bang, bang, bang} can
multiply portrayed impacts; repeated \mention{damn} can strengthen the
conveyed stance. Repetition can thus transform the contribution in
different ways. These examples articulate the contrast to be explained;
they aren't results from a new elicitation study.

The literature motivates those different targets. Potts treats repetition
as potentially strengthening expressive contribution \citep[181--182]{potts2007}.
Dingemanse connects depictive performance with iconic mappings between
form and scene \citep[10--11]{dingemanse2018}. An investigation can
ask whether repetition changes portrayed event structure, conveyed stance,
or both. It needn't assume one answer for every form traditionally
called an interjection.

The descriptive comparison has a definite result. Separate expressive
and depictive groupings preserve a difference in contribution that the
undifferentiated alternative omits. A richer alternative could keep one
superordinate class and record the difference as subtypes or dimensions.
For this descriptive task, that's an acceptable equivalent if it retains
the two relations. The result requires the distinction; it doesn't require
a particular taxonomic rank.

Projective category standing remains a further question. Identifying a use
as depictive from its interpretation under repetition and then predicting
that interpretation from the label would recycle the admission evidence.
An independent test could identify the use from unrepeated contexts and
examine modifications elsewhere. Existing accounts motivate that comparison;
the two illustrations don't establish a general response to repetition or
a complete partition of semantic signs.

The groupings also needn't be disjoint. A depiction can convey an attitude,
and an expressive performance can exploit resemblance. What matters is
whether the relevant contribution is conventional in the selected use and
what follows from it. Semantics can investigate both relations, or several
more, without assigning each to a separate department or requiring one
category to represent the semantic perspective as a whole.

\subsection{Contribution, encoding, and competing analyses}\label{sec:encoding}

Consider the constructed utterance \mention{The damn car won't start},
said after several failed attempts at ignition. The speaker reports a
starting problem and conveys a negative attitude. Calling \mention{damn}
expressive identifies that second contribution. Asking how the wording
helps a hearer recover the attitude raises a different question about
its encoding.

\textcite{wharton2016expressives} develops a procedural account of expressive
meaning, connecting guidance for interpretation with the showing of emotion.
Applied to the car example, such an account would treat \mention{damn}
as a cue for recognizing the attitude being shown. The use could then
be expressive in contribution and procedural in encoding. Those descriptions
would answer different questions about the same utterance.

Now imagine a narrator representing the car's backfire with \mention{Bang!}.
The sign depicts a noise, and its performance can resemble what it
depicts. That iconic presentation differs from \mention{damn}'s conveyed
stance. Identifying the depiction still leaves open how the sign guides
interpretation: resemblance alone doesn't decide whether it supplies
conceptual content, an interpretive procedure, or some combination.

Two accounts would compete if they made exclusive claims
about the same \mention{damn} use. Suppose one assigned it only
conceptual content, such as a negative evaluation, while the other
assigned it only a procedure for recognizing an attitude. These are
hypothetical rivals, not summaries of Potts and Wharton. The first
needn't make annoyance part of what the sentence asserts. Their
incompatible claims about encoding would require adjudication.

Finer groupings also become visible in ordinary situations. Someone
touches a hot pan and says \mention{Ouch!}; someone finds missing keys
under a newspaper and says \mention{Aha!}. The situations are constructed;
\textcite[113]{ameka1992} gives the forms as examples of emotive and
cognitive interjections. The first gives voice to a sensation, the
second displays a change in knowledge or recognition.

Ameka's mental-state grouping is broader than the stance grouping used
here. Finding the keys can be displayed with \mention{Aha!} without
conventionally evaluating them. That reading motivates a cognitive subtype;
it doesn't prevent another occurrence from also conveying affect.
Different contributions, crosscutting accounts of encoding, and rival
encoding hypotheses thus have different consequences even within this
small set of expressions.

\subsection{Interactional resource classes}

Consider two constructed exchanges. A teller says \mention{I tried the
front door}, a listener responds \mention{Mm-hm}, and the teller continues
\mention{but it was locked}. The listener supports continuation of the
story. In a second exchange, someone says \mention{We meet at five};
a recipient who missed the time responds \mention{Huh?}, and the first
speaker repeats \mention{At five}. Here the response prompts repair.

The difference concerns what happens next: the story proceeds after
the continuer, while the meeting time is repeated after the repair
initiator. Closing \mention{bye} has another job, participating in the
end of an encounter. Accounts of interjections in conversation make
these sequential differences central \citep[257--259, 263--265]{dingemanse2024}.
Their union doesn't supply one expectation about the next turn.

For comparison, these can be collected as conventional action formats:
non-elliptical resources that can perform their action without a lexically
specified host. The working class collects resources whose particular sequential
consequences still need separate characterization. Such predictions belong initially to
particular resource classes and environments. The contribution of a
superordinate interactional category would need separate support.

The criterion has consequences at the boundary. Response \mention{okay}
and closing or thanking \mention{thanks} are candidates insofar as their
selected uses conventionally perform the relevant action. A formulaic
greeting \mention{how are you?} can qualify despite having several words.
An ordinary imperative \mention{stop!} doesn't qualify merely because
clause structure permits it to stand alone; a claim that it's also a
conventional routine needs evidence beyond that structure.

Host-dependent \mention{please} in \mention{please sit down} belongs to
a different comparison. A standalone appeal \mention{Please!} would have
to be assessed as that selected use. Thus the admission rule concerns
conventional action formats, rather than brevity, wordhood, or every
expression capable of being uttered alone. These constructed comparisons
specify the proposal; they don't supply an annotated inventory.

The grammatical treatment of \mention{mm-hm} and \mention{huh?} likewise
needs more than standalone occurrence. \textcite[114]{ameka1992} includes the former among
phatic interjections; \textcite[152]{wilkins1992} treats repair \mention{huh?}
as a conventional utterance.
For the selected uses, I adopt the non-elliptical analysis: neither needs
a recovered nominal head or clause to supply its interpretation.
The grammatical assignments below are defeasible analyses supported by
these descriptions, rather than deductions from punctuation.

Routine formulae expose competition within an inquiry.
\textcite[108--110]{ameka1992} separates \mention{bye}, \mention{sorry},
and \mention{thank you} from interjections, whereas
\textcite[142]{wilkins1992} treats routines as a subtype. Both attend
to meaning and interaction. Assigning each a disciplinary label would
misdescribe their disagreement. The selected closing format can be
compared with continuers without settling whether the wider grouping
is the most useful category for either question.

\subsection{Common bearers and distinct questions}

A \term{sign-type} is an identifiable expression in a conventional use.
Its \term{profile} is the combination of grammatical, semantic, interactional,
and other properties attributed to it. A use can retain its identity
while a proposed property turns out to be absent. Otherwise every
surprising observation would identify a different object and leave the
original proposal immune to correction.

The recurrent contrast between standalone and prenominal \mention{damn}
supports distinguishing their conventional uses. In the complement-taking
case, Table~\ref{tab:contrasts} classifies the lexical head \mention{damn}
in that construction. The whole phrase's external distribution is evidence
about its environment, not a replacement bearer. The table records
analytical judgments, not counts or results of a new annotation study.

\begin{table}[htbp]
\centering
\small
\caption{Selected conventional uses and the relations preserved by named analyses.
The columns organize descriptions; they don't allocate one category to each field.}
\label{tab:contrasts}
\begin{tabular}{@{}>{\raggedright\arraybackslash}p{.24\textwidth}
                  >{\raggedright\arraybackslash}p{.31\textwidth}
                  >{\raggedright\arraybackslash}p{.37\textwidth}@{}}
\toprule
Selected bearer & Grammatical analysis & Meaning or action under comparison \\
\midrule
Reaction \mention{ouch!}; \mention{oops!} & Core grammatical interjections & Pain or mishap reaction \\
\addlinespace
Attention-getting \mention{hey!}; silencing \mention{sh!} & Core grammatical interjections & Conative contributions; no evaluative or affective stance required by these selected readings \\
\addlinespace
\mention{gee!}; stance \mention{damn!} & Interjection & Expressive contribution; stance display \\
\addlinespace
\mention{damn} in \mention{damn + NP} & Interjection head or verb: unresolved & Expressive contribution; the head's sequential affordances remain to be specified \\
\addlinespace
Prenominal \mention{damn} & Integrated modifier & Expressive contribution across a grammatical difference \\
\addlinespace
Continuer \mention{mm-hm} & Interjection analysis adopted & Continuation resource; no evaluative stance required by the selected reading \\
\addlinespace
Repair \mention{huh?} & Interjection analysis adopted & Open repair initiator; no evaluative stance required by the selected reading \\
\addlinespace
Closing \mention{bye} & Formulaic standalone expression & Closing resource \\
\addlinespace
Depictive \mention{Bang!} & Interjection or nominal-fragment analysis: unresolved & Sound depiction; stance is a separate question \\
\bottomrule
\end{tabular}
\end{table}

The table records three useful relations. The grammatical core groups
uses with different contributions, including reactions and conatives.
The expressive grouping connects standalone and prenominal \mention{damn}
across a grammatical difference. Within the semantic comparison, stance
and depiction distinguish what the selected \mention{Damn!} and \mention{Bang!}
contribute. A single ordering from less to more interjectional can't
substitute for these inquiry-specific distinctions without losing relevant
information. An ordering may still be useful for a specified task.

\mention{Bang!} makes the unresolved grammatical question concrete.
English also has nominal \mention{bang}, as in the constructed
\mention{a loud bang} and \mention{two bangs}. A standalone occurrence
needs assessment against a nominal-fragment analysis.
Its depictive interpretation can be recognized while that question remains
open. Section~\ref{sec:phonology} completes the related assessment of
what phonological evidence can decide.

\section{What warrants a class, a category, or an explanation?}\label{sec:assessment}

The grammatical core supports expectations connecting independent use,
supplementation, and restricted integration. The expressive grouping preserves
a semantic relation across a grammatical boundary; the expressive--depictive
comparison distinguishes contributions within a semantic inquiry. Assessing
these achievements separately explains how several classifications can
succeed without becoming interchangeable. It also identifies the additional
evidence needed when a descriptive distinction becomes the basis for
a further prediction or a causal explanation.

\subsection{Bearers, grouping, contribution, and alternatives}

Four questions organize the assessment. What bears membership? What groups
the members? What does the grouping help an investigator distinguish,
infer, or explain? And what alternative would serve the same inquiry?
Each proposal also needs a scope: a language or community, relevant
conditions, and, where the claim concerns change, a period. These
questions apply to established proposals and new candidates alike.

The \term{bearer} is the object grouped. Two recordings of \mention{The
damn car won't start} contain two realizations of the selected prenominal
\mention{damn} use. Counting the recordings gives two occurrences, not
two sign-types. A phonologist can classify the individual productions;
a lexical inquiry can ask about the conventional use they realize.
A historian may instead classify a change connecting an earlier sign
with a later one.

The \term{grouping basis} supplies the assignment evidence. Prenominal
\mention{damn} enters the expressive grouping because of its conveyed
stance, while its position inside \mention{the damn car} supplies
grammatical evidence. Grouping it with standalone \mention{Damn!} for
the first reason doesn't erase their difference in the second respect.
Other proposals can use ancestry or participation in a social repertoire;
each needs to identify the relation its membership records.

The \term{contribution} identifies what the grouping accomplishes. A
constructed class can fix a comparison set. A descriptive class can
organize a recurrent relation across independently identified cases.
A category claim requires evidence that co-occurring properties support
induction beyond the application conditions. A causal claim requires
evidence about the process producing an outcome. These achievements
can accumulate, but evidence for an earlier one doesn't establish a
later one automatically.

Descriptive success can be assessed without requiring a prediction about
a differently named property. The expressive grouping preserves a semantic relation across
standalone and integrated constructions; the expressive--depictive contrast
records different contributions concealed by the broad cover term.
Those descriptive achievements establish which similarities an adequate
account needs to preserve. Further projections can build on them,
with their warrant assessed against the evidence for the particular
expectation.

An \term{alternative} can be a broader class, narrower subtypes, a
continuous dimension, or an item-by-item account. In the worked semantic
comparison, separate classes and a superordinate class with explicit
subtypes both preserve the decisive distinction. For an inquiry into
recipient responses, action-specific groupings may be more informative
than their union. The number and rank of categories are outcomes of
comparison, rather than requirements supplied by a field or inquiry.

\begin{table}[htbp]
\centering
\small
\caption{Assessing the expressive grouping without conflating its achievements}
\label{tab:inquiry-record}
\begin{tabular}{@{}>{\raggedright\arraybackslash}p{.17\textwidth}
                  >{\raggedright\arraybackslash}p{.36\textwidth}
                  >{\raggedright\arraybackslash}p{.39\textwidth}@{}}
\toprule
Question & Proposed answer & Evidential consequence \\
\midrule
Bearer & Conventional uses, including standalone and prenominal \mention{damn} &
Identify each use without requiring the target property under investigation. \\
\addlinespace
Grouping & Deployment conventionally conveys evaluative or affective stance &
Distinguish this from describing an emotion, depicting a scene, or managing interaction. \\
\addlinespace
Contribution & Preserve expressive commonality across grammatical constructions &
Distinguish the observed contrast from its extension to new realizations, environments, or sign-types. \\
\addlinespace
Alternative & An undifferentiated interjection class, explicit subtypes, or item-specific accounts &
Compare which semantic relations each preserves and what further generalization each warrants. \\
\bottomrule
\end{tabular}
\end{table}

\subsection{From proposal to category claim}

An admission rule fixes which cases enter a class. A category claim
makes a further commitment: membership supports informative expectations
beyond the evidence used to assign it. Assembling a grouping and
warranting those expectations can succeed or fail separately. Here,
category standing names the latter achievement; it isn't a verdict
on whether a descriptive class is scientifically useful.

Three constructed comparisons distinguish the reach of an expectation.
After identifying the expressive contribution in \mention{The damn car
won't start}, an investigator could examine another speaker's utterance
of that sentence, the question \mention{Has the damn car broken down?},
or \mention{The bloody car won't start}. These comparisons change the
realization, the environment, and the expression respectively.

The question supplies a concrete target: does the negative stance
remain conveyed while the car's breakdown is put in question? Evidence
from declaratives needn't settle that. The investigator can test the
familiar property in a new environment without inventing a differently
named property. Generalizing to independently identified \mention{bloody}
makes a further commitment across sign-types. Neither move is licensed
simply by defining the class to contain the desired cases.

An additional shared property isn't sufficient merely because it's
additional. If all members and nearly every sign outside a proposed
grouping exhibit the target repetition behavior, membership supplies
little additional guidance about it. An imperfect probabilistic association can be
informative when it changes a warranted expectation at the stated
scope. The alternatives in the preceding section help assess that
contribution.

This comparison concerns what membership licenses an investigator to
expect. Whether a category representation improves a particular model
is a separate question, taken up below. Projective warrant doesn't
require that representation to outperform every account of the same
evidence. Both claims remain limited to the profile and scope tested:
support for a subtype doesn't automatically extend to its superordinate
class.

For the grammatical core in Section~\ref{sec:grammatical-core}, the reach
concerns use across environments: evidence about an independently identified
supplementary use supports an expectation of independent utterance.
The cited grammatical descriptions supply the warrant for that connection.
Using both environments to assign a case and then predicting one
from that assignment would instead recycle the identification evidence.

I also adopt the action-specific categories for the sequential relations
documented in the cited descriptions. The same evidential policy applies
across inquiries: established descriptions can warrant scoped semantic and
interactional projections without a new experiment in this paper.

The semantic comparisons here chiefly establish how expressive commonality
crosses grammatical distinctions and how stance differs from depiction.
Their focus reflects the question being pursued. Further expressive
generalizations can draw on existing evidence just as the grammatical
generalization does; each is assessed at the scope that evidence supports.

\subsection{Projection and identity}

Revising the analysis of a conventional use can improve subsequent
predictions without turning an earlier failure into a success. To
keep those judgments separate, a projective claim states what's been
observed, what's expected, and over which bearers and conditions.
Its \term{warrant} is the evidence supporting that expectation
\citep{Goodman1955}. Claim content, evidential status, and proposed causal
ground can then be assessed separately
\citep[11--12]{reynolds2026projectibilityHistory}.

Suppose source contexts identify an expressive use and the target is
its interpretation under repetition. Repetition can't also supply the
case-level evidence used to admit that use. Similarly, if resistance
to determiner combination is the target, an unexpected determiner
can't by itself justify excluding the case as conversion. A change
to the identification rule creates a revised claim requiring assessment.

Imagine an investigator who's identified expressive \mention{damn}
only in standalone reactions such as \mention{Damn! We've missed the
train}. Their provisional identification uses the reaction meaning;
it doesn't require standalone position. They predict that these uses
resist grammatical integration. Encountering \mention{The damn car won't
start} then presents an adverse case: the expression conveys stance
while functioning as an integrated modifier inside the noun phrase.

This is a hypothetical research sequence using familiar constructions.
The investigator should check the structural analysis, but the confirmed
prenominal occurrence counts against the restriction under the original
identification rule. Its inconvenient distribution isn't by itself a
reason to remove it from that assessment. Later recognizing a distinct
prenominal use can improve the analysis without making the original
prediction successful.

Evaluate the original projection under the identification rule available
when it was made. Later evidence may warrant distinguishing a recurrent
conventional use by its contribution or distribution. That revision
can improve subsequent predictions, but it doesn't remove earlier
counterexamples from the original assessment. Observations that motivate
a new distinction don't also supply independent confirmation of its
predictive success.

A quoted token or a productive noun denoting an utterance remains
another possible analysis, but it needs evidence for quotation or
nominal reference. A correction to a mistaken structural analysis should
be recorded as such. It differs from changing the identification rule
so that a confirmed counterexample no longer falls within its scope.

An informative restriction also needs relevant opportunities. Failure
to observe inflection says little if the sampled environments wouldn't
elicit it under a rival analysis. A grammatical test needs source and
target constructions; an interactional test needs actions and sequential
environments. The common bearer allows these claims to concern the
same use. It doesn't make their evidence interchangeable.

\subsection{Representation and causal explanation}

Suppose an inventory records that a use occurs standalone and conveys
stance. Adding the label \enquote{standalone stance use} abbreviates
those two observations; it adds no information to a fully observed
record. The combined grouping could still be useful when investigating
something not yet recorded, such as interpretation under repetition.
Pooling observations from its members might help estimate that pattern
from sparse records. Whether it does is a comparative question.

The grouping selects an intersection: stance \mention{Damn!} qualifies,
while prenominal \mention{damn} in \mention{the damn car} fails the
standalone condition. A model can represent that intersection as one
label or keep its two component assignments. Their relative value
depends on the information available and the target being predicted.
A benefit for estimating repetition behavior wouldn't establish a
unique cause of the shared properties or priority for every inquiry.

Explanation asks another question. Once the present form of \mention{gee}
has been recorded, knowing its religious source adds nothing to that
already observed spelling. Its history can still explain how the form
arose. More generally, conditioning a prediction on consequences of a
historical process can remove the association through which that process
is recognized. Section~\ref{sec:history} distinguishes this predictive
comparison from assessment of the proposed history itself.

\textcite[295--299]{hill1965} offered viewpoints for comparing causal
interpretations with alternatives, including temporal order and consistency
with independently supported knowledge. These concern specified causes
and outcomes. In the history of \mention{gee}, for example, the
sequence of semantic change and clipping helps distinguish the processes
that contributed to the present form (Section~\ref{sec:history}).

\section{Phonological organization and situated realization}\label{sec:phonology}

Shared phonological possibilities can justify comparing signs across
grammatical categories. The grammatical uncertainty about depictive \mention{Bang!}
therefore leaves a phonological comparison available. Such comparisons
can concern sounds or combinations unusual elsewhere in a variety,
or the prosodic packaging of particular occurrences.

These observations motivate different classifications. The first can group
sign-types by their phonological possibilities. The second groups realizations
by how speakers produce them in context. Their relation connects
the resources available for a conventional sign with the ways speakers
realize it on particular occasions.

\subsection{A phonological stratum of signs}

English conversational clicks make the first possibility concrete. The
click represented by \mention{tut-tut} differs from an ordinary spelling
pronunciation \citep[147]{wilkins1992}. Forms such as \mention{ugh} permit
sounds outside the usual phoneme inventory of standard British and
American English \citep[853]{quirk1985}. A phonological inquiry can
ask which conventional signs admit such resources, how their variants are
organized, and how speakers acquire or adapt them.

\textcite[112--113]{ameka1992} diagnoses a specific inference: unusual
phonology shared with interjections has been used to include onomatopoeic
words in the grammatical interjection category. His alternative is to
recognize a phonological grouping whose members can belong to different
grammatical categories. The following assessment develops the evidential
consequence of that distinction.

The grammatical claim concerns membership in a lexical category. The
evidence offered concerns sounds or sound patterns unusual elsewhere in
the language. That evidence supports a comparison of phonological
possibilities across independently identified signs. By itself, it doesn't
settle whether a particular use has nominal, verbal, or interjectional
grammar. Nor does sound depiction itself entail phonological anomaly: resemblance
and distribution within the sound system are different relations.

Phonological evidence can support a grammatical expectation where an
independently established association warrants it. The shared unusualness
considered here doesn't itself establish that association. The limitation
concerns the evidence connecting the properties, not the fact that they
belong to different domains.

For English \mention{Bang!}, the competing grammatical analysis is a
standalone nominal fragment, supported as a possibility by the independently
available noun in \mention{a loud bang}. A depictive interpretation leaves
that analysis open. Even if an occurrence had unusual phonation, the
additional observation would concern its realization. Choosing between
nominal and interjectional analyses needs distributional and constructional
evidence about the selected conventional use, or an independently warranted
association connecting the observed phonology to that grammar.

The strongest alternative retains one broad descriptive class and records
phonology and grammar as separate attributes. For organizing known cases,
that can preserve the same information. It becomes inadequate only if
membership in the broad class is used to infer a grammatical assignment
without the relevant evidence. The assessment rejects
that inference while retaining the cross-grammatical phonological
comparison. It doesn't establish a single projectible phonological stratum.

An account of the stratum would gain projective reach if membership
organized acquisition, adaptation, or alternation beyond the diagnostic
used to identify it. Consider the contrast between producing a click
and pronouncing \mention{tut-tut} as spelled words. An investigation
could first group signs by independently recorded articulations, then
ask which permit a spelling pronunciation and in what settings.
That prospective comparison would test a relation beyond the articulation
used for grouping; the reported alternatives don't yet establish it.

Regular phonology matters too. In a ten-language study of open repair
interjections, \textcite{dingemanse2013huh} find broad phonetic similarities
alongside language-specific vowel targets and intonation. The repair interjections in Icelandic and
Cha'palaa have falling contours, consistent with local interrogative
prosody. The study concerns comparable repair environments. It supplies
no reason to assign every occurrence resembling \mention{huh} to one
phonologically exceptional class.

Phonological classification thus has work to do within ordinary linguistic
organization as well as at its margins. An inquiry into repair intonation
could group forms with interrogative resources; another into unusual
segments could group conversational clicks with other marked forms.
Both could be useful. A single label such as \term{phonological interjection} would obscure
the difference unless the inquiry and grouping basis were specified.

\subsection{Prosodic classes of realizations}

A prosodic boundary is a property of a realization. A sign-type can
characteristically occupy its own intonation phrase while allowing different
realizations in different environments. Conversely, realizations of
grammatically different signs can share a boundary pattern. Supplements
provide a familiar example: interjections, noun phrases, and relative
clauses can all be prosodically detached
\citep[66, 1350]{HuddlestonPullum2002}.

Take a recording of \mention{Damn, we're going to be late!}. The
context can identify the reaction use while the acoustic record supplies
evidence about a boundary, such as a pause or pitch movement. A
second recording can realize the same use with different timing.
Classifying these productions asks how that occurrence was packaged;
a claim about characteristic packaging needs the distribution of
productions across relevant environments.

Interaction supplies some of those environments. \textcite{ogden2013clicks} describes English
clicks and percussives in stance-taking, sequence management, word searches,
and turn timing, with visual behaviour contributing to their interpretation. A classification of productions by timing and
articulation can intersect an analysis of sequential practices.
The range of jobs also cautions against assigning one conventional action
or expressive meaning to every click.

The comma in a transcript of \mention{Damn, we're going to be late!}
is a different kind of observation. It may record a transcriber's
analysis of detachment, but it isn't itself a pause or a grammatical
relation. Entering that comma once as a prosodic boundary and again
as syntactic independence would manufacture apparent agreement between
the measures. An acoustic analysis and a constructional analysis can
converge without relying on that single transcription decision.

The two phonological inquiries extend the family in different ways. Phonological strata can share their sign-type bearers and
enter membership intersections. Prosodic classes of occurrences connect
to them through realization and through distributions over contexts.
Treating both as extra columns in the same sign-type table would hide
the relation that the phonological inquiry is meant to explain.

\section{Historical ancestry and classes of change}\label{sec:history}

The history of \mention{gee} explains a feature of its present form
that a synchronic interjection label leaves unexplained. Signs with
different current distributions can descend from related sources, while
signs with similar current profiles can arrive by different routes.

A historical inquiry can classify present signs by independently established
ancestry. It can also classify changes, such as clipping or reanalysis.
The first has sign-types as bearers; the second has changes linking
sources and outcomes.

\subsection{The documented route to \mentionhead{gee}}

\textcite{gehweiler2008} traces \mention{gee} through an invocation of \mention{Jesus},
an exclamatory use of the religious name, and a clipped form. In the invocation, the speaker addresses the deity.
The exclamation can express a reaction without that address. Clipping
then produces a form whose relation to the name can become opaque.
The sequence separates semantic change from the later change of form.

The transitions also have different proposed explanations. Gehweiler
relates the first to gradual invited inference, aided by the syntactically
peripheral invocation. She analyzes the clipping as an abrupt, taboo-motivated
change involving hearer reanalysis \citep[84--88]{gehweiler2008}. The
source invocation already permits peripheral packaging.

The history explains why the shape of \mention{gee} bears a relation
to religious invocation even for a speaker who doesn't recognize that
relation. Its present opacity alone wouldn't identify its ancestry.
The documented sequence supplies the additional evidence for a
restricted instance of path dependence: a particular route helps explain
a feature of its outcome, without making every other outcome of that
route predictable.

The explanation concerns a particular outcome. It doesn't yet establish
that a class of signs with religious ancestry outperforms alternatives
or supports a common projection. That comparison would require independently
documented histories and outcomes across more than this one route.

The distinction also corrects an ambiguity in \term{primary} and
\term{secondary} interjections. Ameka's distinction concerns whether
a form has other, non-interjectional uses in the language
\citep[105--106]{ameka1992}. It isn't a partition into historically
underived and historically derived signs. Gehweiler's primary
\mention{gee} has a documented derived history. Conversely, an unknown
etymology supplies a gap in knowledge, not a common origin shared by
all forms whose sources remain unknown.

\subsection{Source, mechanism, transmission, and trajectory}

A useful historical description separates four relations. The
\term{source} identifies an earlier sign or construction. The
\term{mechanism} specifies a change, such as clipping or reanalysis.
The \term{transmission setting} concerns how the innovation spreads,
including borrowing where supported. The subsequent \term{trajectory}
concerns further changes in use or form. Signs can share one relation
without sharing the other three.

For \mention{gee}, a religious invocation supplies the source and
tabooistic clipping supplies one mechanism. Those facts don't establish
the transmission history of every later community's use. Nor does a
shared source force a common trajectory: an exclamatory religious name
can remain phonologically close to its source without later undergoing
clipping. Historical classes can be organized at several
levels, according to which relation the inquiry needs.

The proposed verbal route to complement-taking \mention{damn} illustrates
unequal support among histories. \mention{God damn him} has the shape
of an optative, a construction independently recognized in English
\citep[944]{HuddlestonPullum2002}. Omission of the subject and increasing
conventional expressive force offer a plausible route to \mention{damn
him}. The surviving complement fits that reconstruction. These synchronic
comparisons don't establish the sequence as independently as Gehweiler's
historical record does.

The source optative already uses plain-form \mention{damn}, so the
reconstruction doesn't require loss of tense marking at that transition.
And the present complement remains compatible with competing analyses
of the head as verbal or interjectional. A classification by independently
established verbal ancestry could cross that present grammatical boundary.
Inferring the ancestry solely from the complement and then invoking
the ancestry to predict that same complement would be circular.

Classes of changes require a different representation. An instance of
clipping relates an earlier form to an outcome; an instance of semantic
reanalysis relates earlier and later interpretations. Neither is itself
a sign-type alongside \mention{gee} and \mention{damn}. A process
classification can support comparisons among developments, while explicit
source--result relations connect those developments to classifications
of their inputs and outputs.

\subsection{What historical evidence can establish}

Historical explanation and forward projection need different comparisons.
A documented route may explain a present feature through the changes
that produced it. A forward projection asks what information about the
route warrants expecting in another case. For a verbal-source class,
that target could concern surviving restrictions on dependents, provided
those restrictions weren't used to assign the source. The present
examples motivate that inquiry without establishing a category-wide result.

Evidence for replication also needs to distinguish innovations from their
descendants. Suppose one community coins a clipped exclamation and
three others borrow it. Four community records then document one
origin, not four independent instances of clipping. They can support
an account of transmission or subsequent divergence. Testing whether
the same source repeatedly gives rise to the same outcome would
require independently established innovations. This hypothetical count
makes the sampling unit depend on the historical claim.

An ancestry classification can still explain a present contrast even
when ancestry adds no predictive information after all its consequences
are observed. Conditioning on those consequences changes the question
to whether ancestry contributes anything further. That can be a useful
test of predictive economy; it isn't a general test of whether the
documented process occurred or helped produce the contrast.

Historical classes can have two kinds of payoff. They can
organize independently supported source relations among present signs,
and they can compare processes across developments. These are legitimate
descriptive and explanatory contributions. An empirical ancestry category,
in the stronger usage adopted here, also needs support for the
inductions it licenses. The documented history of one sign can't supply
that support for an entire source class.

\section{Broader classes and further inquiries}\label{sec:extensions}

The traditional interjection domain also intersects classes whose most
informative members lie elsewhere. Social evaluation can group forms
by their appropriateness for audiences and settings. Processing inquiries
can group familiar expressions by how they're produced or understood.
An interjection's membership in either class needn't make the class
interjection-specific. Three questions remain distinct: whether the
broader class is supported, how it intersects the traditional domain,
and whether an \term{interjection} label helps describe that intersection.

\subsection{Social evaluation and recognized repertoires}

Religious exclamations connect the historical case to a further classification.
In Gehweiler's account, taboo pressure contributes to clipping \mention{Jesus}
into \mention{gee}. A present-day inquiry can group uses by their
evaluated appropriateness for an audience or setting. It asks how
speakers assess the available forms in their own community.

\textcite{jayjanschewitz2008} examine ratings of hypothetical swearing scenarios,
finding that judgments depend on the word, speaker--listener relations,
and social and physical context. These ratings support contextual
differentiation in evaluations of taboo language.

Consider a proposed comparison of \mention{Gee!} and \mention{Jesus!}
as reactions to surprising news among close friends and during a
formal interview. Ratings would ask how appropriate each choice seems
in each setting. A difference in those ratings would distinguish
social appropriateness without by itself changing the grammatical classification
of the reaction.

Social indexing poses a related but distinct inquiry. Agha's account
of enregisterment concerns the social recognition of speech forms as
indexing attributes of speakers, with a repertoire's associations and
social reach subject to change \citep[38--40]{agha2005}. A classification
by participation in a recognized repertoire could cross grammatical,
semantic, and sequential divisions. The relevant evidence concerns
recognition and use of those associations, beyond a frequency difference
between geographical collections.

For example, if listeners regularly associated \mention{Gee!} with an
old-fashioned speaker persona, their judgments would support a social
association. They wouldn't merely be rating whether it was acceptable
in the interview. Both investigations could concern the same surprise
reaction, while asking different questions about its social meaning.
Such associations could also unite forms used for several actions.
Any claimed convention would need a specified community and evidence
of shared recognition.

\subsection{Formulaicity and processing}

Compare a familiar greeting, \mention{How are you?}, a pause filler,
\mention{um}, and the idiom \mention{kick the bucket}. These illustrative
forms differ in grammar and conversational job. A processing inquiry
could ask whether speakers produce them as familiar routines.
Its grouping would cross the earlier grammatical and semantic divisions;
\term{formulaic expressions} might name it more clearly than another
kind of interjection.

\textcite{vanlanckersidtis2015} report differences in formulaic
language production and comprehension across clinical groups and
healthy participants. Their production
analysis combines six kinds of formulaic material, including conversational
formulae, pause fillers, sentence stems, and idioms. The results give
empirical consequences to a grouping that crosses grammatical categories:
pause fillers can be studied alongside phrases and idioms despite
their grammatical differences.

For the greeting, one question is whether \mention{How are you?} is
retrieved fluently because repeated use has helped establish a routine.
A comparison at the same measured frequency asks how fluency differs
among expressions with comparable recorded use. That adjustment doesn't
by itself identify how much repetition contributed to the greeting's
fluency. A causal interpretation still needs assumptions about exposure,
retrieval, and how the frequency measure relates to speakers' experience.

Suppose retrieval across \mention{How are you?}, \mention{um}, and
\mention{kick the bucket} tracked familiarity more closely than any
interjection grouping. That result would favour the broader or continuous
representation for this processing inquiry. It wouldn't give the filler
\mention{um} the internal grammar of \mention{kick the bucket}. The
grammatical contrast would remain available for a different question,
provided its own evidence held.

\subsection{Further inquiries and their objects}

Morphological analysis can distinguish a retained plural shape in
exclamatory \mention{heavens!} from a productive number contrast. Attested
noun and verb uses of \mention{ouch} in child-directed speech likewise
make the conventional use consequential for grammatical classification
\citep[36]{stange2016}. The lexical profile of a reaction doesn't describe
every construction containing the same spelling.

Other inquiries change the object under investigation. A child's acquisition
of \mention{ouch} concerns a learner and a series of learning episodes.
A recording of \mention{Bye!} with a wave includes visible timing absent
from a lexical inventory. In a constructed comic panel, \mention{BANG}
can depict a noise without representing anyone's speech. Learning,
multimodal performance, and graphic depiction connect to conventional signs
through different relations, each requiring its own comparison.

Annotation supplies a practical example. A tag can retrieve conversational
responses while combining several grammatical or semantic classes. Its
success concerns retrieval under a coding policy. To use the retrieved
tokens as evidence about a lexical category, an investigator needs
to establish how that policy relates to grammatical membership. The
following measurement case shows the consequence of leaving that relation
unchecked.

\section{A measurement case: tags, strings, and uses}\label{sec:evidence}

A retrieval class can be useful for a corpus query and still give
a misleading picture when treated as a linguistic category. In the
Corpus of Global Web-Based English (GloWbE), which contains country-labelled
web sections \citep{daviesfuchs2015}, the spelling \mention{ha} retrieves
both reactions and abbreviations for hectares. A tag, a string, and
a conventional use select different sets. The distinction changes the
apparent geographical pattern.

The constructed strings \mention{12 ha of land} and \mention{Ha! I
knew it} make the difference visible. A search for the spelling
\mention{ha}, ignoring capitalization, can retrieve both. An investigation
of the reaction use needs to distinguish them. A tag intended to
make that distinction is still a measurement decision whose accuracy
needs checking.

The \mention{ha} record makes the consequence concrete. Some tagged
occurrences abbreviate \term{hectare}, especially in agricultural prose.
A rule-based screen of exported contexts changes Kenya's rank from fifth
to fifteenth or sixteenth under six treatments of laughter and unclear
cases; pooling gives ranks from fourteenth to sixteenth. The large
movement is stable across those choices. Its exact size remains
conditional on the screen and on the exported samples.

The screen was tuned against forty Kenyan and Tanzanian examples;
its resulting 39/40 binary agreement is a development check. It doesn't
establish equivalent accuracy in every country section. The observed
geographic distribution depends on both the linguistic
phenomenon and the measurement procedure. An analysis of category
variation would need independently identified uses and comparable
opportunities for their occurrence.

Similarly, changing frequency ranks in the Corpus of Historical American
English (COHA) describe a tagged inventory \citep{davies2010}.
They don't establish persistence of a semantic or grammatical profile.
Two further analyses, of derivatives and \mention{fie}, fail their
intended confirmatory interpretation: the primary derivative pipeline
assigned semantic residue rather than measuring it, and the \mention{fie}
input retained ineligible optical character recognition substitutions.
Their confirmatory claims are withdrawn.

The \href{inquiry-relative-evidence-supplement.pdf}{evidence supplement}
preserves the complete records: GloWbE measurement and model comparisons
(Section~A), the derivative audit (Section~B), the \mention{fie} audit
(Section~C), and inventory turnover (Section~D). Its retained model
comparison predicts an omitted country--item cell within the observed
design; it doesn't assess new countries, new items, or the proposed
linguistic groupings. These records locate
the failure at the relevant claim. A useful retrieval class survives
as a tool, but its output can't establish a linguistic category
without evidence connecting that class to independently identified uses.

\section{The general consequence}\label{sec:discussion}

A classification can preserve one supported pattern of similarities
while crossing another. Conative \mention{sh} shares the grammatical profile
of \mention{ouch} without requiring its affective contribution. Prenominal
\mention{damn} shares expressive properties with standalone \mention{Damn!}
without becoming an intermediate grammatical interjection. The relevant
relations are available for different inquiries; failure to preserve
them all in one grouping needn't be a classificatory defect.

Depictive \mention{Bang!} introduces a further semantic relation, and continuers
and repair initiators organize different sequential environments. These
distinctions prevent pluralism from becoming a fixed arrangement with
one category for grammar, another for semantics, and a third for interaction.
The number and extension of the groupings depend on the questions
and the relations supported by the evidence.

\subsection{Many-to-many relations and local failure}

Dupr\'e's pluralism and Khalidi's crosscutting kinds make room for
multiple defensible groupings. The linguistic cases add a practical
requirement: specify which relation a proposed grouping preserves and
which objects participate in it. Haspelmath's division between descriptive
categories and comparative concepts supplies a related lesson about
analytical purpose. Here, the same care is needed within a language
and within an inquiry.

In \mention{The damn car won't start}, expressive contribution concerns
the attitude conveyed; procedural encoding concerns how the wording
guides its recognition. Both descriptions could hold. The contrast with
depictive \mention{Bang!} concerns a different contribution, while the
exclusive encoding hypotheses in Section~\ref{sec:encoding} disagree about
the same use. Assigning the hypotheses to different fields wouldn't
make their claims compatible.

Failure has correspondingly local consequences. Suppose a repetition
test found that an independently identified expressive use didn't
strengthen stance under the specified conditions. The prediction of
intensification would fail there; the observed distinction between
conveying stance and depicting a sound wouldn't disappear with it.
Restricting the prediction to another subtype would create a new
claim. Its support couldn't be inherited from the failed generalization.

The resulting pluralism combines positive classifications with decisions
about particular claims. The grammatical core supports distributional
expectations; the semantic comparisons preserve expressive and depictive
relations that a grammatical inventory alone leaves unspecified.
The phonological comparison rejects an unsupported transfer to grammatical
membership, while the \mention{gee} history supports a particular explanation.
The success of one analysis gives no general priority over the others.

\subsection{Overlap, mechanisms, and category standing}

Selecting standalone stance uses includes expressive \mention{Damn!}
and excludes prenominal \mention{damn} in \mention{the damn car}. That
intersection is well-defined and can fix a useful comparison set.
Its further value depends on the question: does restricting attention
to those uses help with the proposed inference? Neither the overlap
itself nor the availability of its component assignments answers that.

Distinct categories needn't be statistically independent, and they can
share causal processes. Strong correlation doesn't make two questions
identical; a shared history doesn't establish every present property.

For \mention{Bang!}, uncertainty between a nominal-fragment and an
interjection analysis doesn't itself establish graded membership. The
analysis may simply be unresolved. A separate inquiry could ask how
typical speakers judge it compared with \mention{Damn!}; prototype
accounts organize such typicality and similarity \citep{rosch1975,taylor2003}.
Either a graded or a categorical representation can serve a task,
but its form alone doesn't establish what it licenses an investigator
to infer.

A field supplies no fixed quota of categories, and an inquiry
needn't end with just one. A grouping may serve several questions across
fields, but warrant only specified projections under specified conditions.
The historical label \term{interjection} identifies a meeting point for
these inquiries. The evidence, rather than the label or departmental
boundary, determines which proposed categories deserve recognition.

\section{Conclusion}\label{sec:conclusion}

English interjections support several classifications because the relevant
similarities crosscut. A grammatical core including \mention{ouch},
\mention{oops}, \mention{hey}, and \mention{sh} supports expectations about
independent utterance and clausal supplementation. Within that core,
reactions and conatives differ in contribution. Across the grammatical boundary,
standalone and prenominal \mention{damn} retain expressive commonality.
The semantic comparison with depictive \mention{Bang!} adds another
distinction within the same broad inquiry into meaning.

Crosscutting classifications needn't be rival approximations to a single
complete classification. Each can preserve a supported relation while
disregarding another. Recognizing those relations separately changes the
diagnosis of apparent anomalies: a conative use can be a grammatical
interjection without expressing evaluative stance, and an integrated modifier
can share expressive properties with an interjection without sharing
its lexical category.

The same approach distinguishes coexistence from competition. Expressive
commonality leaves the rival grammatical analyses of complement-taking
\mention{damn} unresolved; it doesn't reconcile them. A later distinction
between conventional uses can improve predictions without erasing an
earlier failure. And a retrieval class can remain useful after its
output has proved misleading as a measure of linguistic uses.
Pluralism permits several classifications to succeed while keeping particular
claims answerable to evidence.

\ifblind\else
\section*{Acknowledgements}

This paper reanalyses material I compiled for the Wikipedia article
\enquote{English interjections}.
\href{https://github.com/BrettRey/interjections-hpc/tree/master}{The archived
data and analysis code} accompany the linked evidence supplement.
\fi

\clearpage
\begingroup
\setlength{\emergencystretch}{3em}
\printbibliography
\endgroup
\end{document}
